\chapter{Notes}
\label{ch:notes}

% ======================================================================
\section{Units}
\label{sec:units}
% ======================================================================


% ----------------------------------------------------------------------
\subsection{Phonopy / QE Units}
\label{sec:phonopy_units}
% ----------------------------------------------------------------------

Phonopy stores force constants in eV/\AA$^2$, masses in amu, and
positions in \AA.  When phonopy is used with QE, the internal representation
uses atomic units (bohr, Ry/bohr$^2$).  Everything is normalised to the
eV/\AA\ system upon loading.

The dynamical matrix eigenvalues $\lambda$ have units
eV/(amu$\cdot$\AA$^2$).  Phonopy converts to THz via
\begin{equation}\label{eq:VaspToTHz}
  f\;[\text{THz}]
  = \mathrm{sgn}(\lambda)\;\sqrt{|\lambda|}
    \times \underbrace{
      \frac{1}{2\pi}\sqrt{\frac{e}{m_u \cdot A^2}}
    }_{\displaystyle \approx 15.633\;\text{THz}},
\end{equation}
where $e = 1.602\times10^{-19}$\,J, $m_u = 1.661\times10^{-27}$\,kg,
$A = 10^{-10}$\,m.

When loading a phonopy calculation performed via the QE interface,
lengths are in bohr and force constants in Ry/bohr$^2$. Conversion is
\begin{align}
  \text{lattice [}\text{\AA}\text{]}
    &= \text{lattice [bohr]} \times a_0,
    &a_0 &= 0.52918\;\text{\AA/bohr}, \\
  \Phi\;[\text{eV/\AA}^2]
    &= \Phi\;[\text{Ry/bohr}^2] \times \frac{E_h/2}{a_0^2},
    &\frac{E_h/2}{a_0^2} &\approx 48.59\;\text{eV\,bohr}^2\text{/(Ry\,\AA}^2).
\end{align}


% ----------------------------------------------------------------------
\subsection{quatrex Units}
\label{sec:solver_units}
% ----------------------------------------------------------------------

The quatrex NEGF solver operates in angular frequency squared:
$\omega^2$ in $(\mathrm{rad/s})^2$.  The conversion factor from
phonopy eigenvalue units eV/(amu$\cdot$\AA$^2$) to $(\mathrm{rad/s})^2$
is
\begin{equation}\label{eq:conversion}
  C = \frac{e}{m_u \cdot A^2}
  \approx 9.648 \times 10^{27}\;\text{(rad/s)}^2.
\end{equation}
The real-space blocks stored for the solver are therefore
$H(\vc{R}) = C \cdot \Phi(0,\vc{R}) / \sqrt{m\,m'}$, in units of
$(\mathrm{rad/s})^2$.


\section{Quatrex input format}
\label{sec:quatrex_input}

The quatrex NEGF solver reads three input files that together specify
the real-space Hamiltonian, the crystal structure, and the simulation
parameters.


% ----------------------------------------------------------------------
\subsection{Input files}
\label{sec:input_files}
% ----------------------------------------------------------------------

\subsubsection{dynamical\_matrix.mat}

A MATLAB \texttt{.mat} file containing real-space blocks $H(\vc{R})$
in Convention~B (see \cref{sec:dynmat}) (i.e. Force Constants), in units of $(\mathrm{rad/s})^2$.  Each key is a
string \texttt{[nx, ny, nz]} mapping to an
$(n_\mathrm{dof} \times n_\mathrm{dof})$ real matrix, where
$n_\mathrm{dof} = 3N_b$ and $N_b$ is the number of atoms per unit
cell.  For nearest-neighbour interactions, the file contains up to 27
blocks with $n_x, n_y, n_z \in \{-1, 0, +1\}$.

The solver loads these blocks, constructs a BTD matrix along the
transport direction, and performs the transverse Fourier sum
\begin{equation}
  D(\vc{k}_\perp)
  = \sum_{(n_x, n_y)} M(n_x, n_y)\;
    e^{2\pi i\,(k_x n_x + k_y n_y)},
\end{equation}

\subsubsection{structure.xyz}

An extended XYZ file with lattice vectors in the header (ASE
convention):
\begin{verbatim}
8
Lattice="5.4018 0.0 0.0 0.0 5.4018 0.0 0.0 0.0 5.4018" ...
Si  0.00000000  0.00000000  0.00000000
...
\end{verbatim}
The loader (\texttt{distributed\_read\_xyz}) extracts the $3\times3$
lattice matrix and Cartesian atom positions.  The number of atoms times
\texttt{num\_orbitals\_per\_atom} determines $n_\mathrm{dof}$.

\subsubsection{quatrex\_config.toml}
See \texttt{src/quatrex/core/config.py}
\begin{center}
\begin{tabular}{lll}
  \toprule
  Field & Example value & Purpose \\
  \midrule
  \texttt{transport\_direction}       & \texttt{'z'}              & BTD direction \\
  \texttt{construct\_from\_unit\_cell}& \texttt{true}             & input mode (see below) \\
  \texttt{num\_transport\_cells}      & 4                         & device length in unit cells \\
  \texttt{neighbor\_cell\_cutoff}     & \texttt{[1,1,1]}          & maximum block range \\
  \texttt{kpoint\_grid}               & \texttt{[4,4,1]}          & transverse MP mesh \\
  \texttt{kpoint\_shift}              & \texttt{[0,0,0]}          & mesh offset \\
  \texttt{num\_orbitals\_per\_atom.Si}& 3                         & DOFs per atom \\
  \texttt{phonon.eta}                 & \texttt{1e-8}             & broadening $(\mathrm{rad/s})^2$ \\
  \texttt{phonon.obc.algorithm}       & \texttt{"sancho-rubio"}   & OBC method \\
  \bottomrule
\end{tabular}
\end{center}

% ----------------------------------------------------------------------
\subsection{Input modes}
\label{sec:input_modes}
% ----------------------------------------------------------------------

quatrex supports two modes selected by
\texttt{construct\_from\_unit\_cell}.  In both modes the input is a
MATLAB \texttt{.mat} file loaded by \texttt{scipy.io.loadmat}.  The
loader (\texttt{distributed\_load} in
\texttt{qttools/utils/mpi\_utils.py}) parses keys that start with
\texttt{[} and interprets them as integer tuples:
\begin{lstlisting}
arr = loadmat(path)
arr = {
    tuple(map(int, r.strip("[]").split(","))): h_r
    for r, h_r in arr.items()
    if r.startswith("[")
}
\end{lstlisting}
i.e., a key string \texttt{"[0, 1, -1]"} becomes the Python
tuple \texttt{(0, 1, -1)}.  Each value is a 2D NumPy array (real or
complex) representing a matrix block.

\subsubsection{Mode 1: Pre-assembled device matrix
  (\texttt{construct\_from\_unit\_cell = false})}

The \texttt{.mat} file contains a single pre-assembled device matrix
(already tiled and including all transport cells), stored under the
key \texttt{(0,0,0)}:
\begin{lstlisting}
matrix_sparray = distributed_load(f"{matrix_name}.mat")
assert (0, 0, 0) in matrix_sparray.keys()
matrix_sparray = matrix_sparray[(0, 0, 0)]
\end{lstlisting}
The result is a single sparse matrix of size
$(N_\mathrm{blocks} \times B)^2$, where $B$ is the block size and
$N_\mathrm{blocks}$ is the number of transport cells.

No transverse $k$-point Fourier assembly is performed, since
the matrix is pre-assembled and has no dictionary of periodic shifts.
This mode is suited for non-periodic heterostructures.

\subsubsection{Mode 2: Unit-cell blocks
  (\texttt{construct\_from\_unit\_cell = true})}
\label{sec:mode2}

The \texttt{.mat} file contains real-space blocks $H(n_x, n_y, n_z)$
of the unit cell, keyed by lattice translation indices.

\paragraph{Loading and grid construction.}
The function \texttt{\_load\_matrix\_from\_unit\_cell} loads all keys,
determines the bounding box from the minimum and maximum key
coordinates, and packs the blocks into a 5D array:
\begin{lstlisting}
# keys: array of all (nx, ny, nz) tuples
min_coords = keys.min(axis=0)
max_coords = keys.max(axis=0)
grid_shape = max_coords - min_coords + 1
unit_cells = zeros(grid_shape + matrix_shape)
for coord, matrix in matrices.items():
    unit_cells[coord] = matrix
\end{lstlisting}
The resulting array has shape
$(G_x, G_y, G_z, n_\mathrm{dof}, n_\mathrm{dof})$, where
$(G_x, G_y, G_z)$ is the grid shape (e.g.\ $(3, 3, 3)$ for
$\{-1,0,+1\}^3$).

The loader checks that
$G_x \times G_y \times G_z = \text{number of keys}$, i.e.\ the grid
must be \emph{dense}: no missing entries are permitted.  Missing
blocks should be stored as zero matrices.

\paragraph{Neighbour cell cutoff.}
If \texttt{neighbor\_cell\_cutoff} is set (e.g.\ \texttt{[1,1,1]}),
\\\texttt{trim\_tight\_binding\_matrix} zeros out all blocks beyond that
range.  For cutoff $(c_x, c_y, c_z)$, a block at index
$(n_x, n_y, n_z)$ (centred at zero) is kept only if
$|n_x| \le c_x$, $|n_y| \le c_y$, $|n_z| \le c_z$.

\paragraph{Transport cell (supercell) construction.}
The function \texttt{expand\_tight\_binding\_matrix} converts the
unit-cell blocks into a BTD sparse matrix for a given transverse
periodic shift.  Along the transport direction $z$:
\begin{enumerate}
  \item The supercell size in the transport direction (i.e. number of uc in
        transport direction) is $G_z / 2$ (integer division), and 1 in the transverse
        directions.  For $G_z = 3$, this gives a supercell of one unit
        cell ($3 // 2 = 1$).
  \item Three BTD blocks are extracted for each transverse periodic
        shift $(n_x, n_y)$:
        \begin{itemize}
          \item $n_z = -1$: lower off-diagonal (coupling from the
                previous layer),
          \item $n_z =  0$: diagonal (on-site),
          \item $n_z = +1$: upper off-diagonal (coupling to the next
                layer).
        \end{itemize}
  \item Each block is assembled by \texttt{\_get\_transport\_block},
        which loops over all pairs of local shifts within the supercell
        and selects the appropriate unit-cell block for each pair.
  \item The three blocks are tiled $N_\text{transport}$ times along the
        diagonal to build the full BTD sparse matrix.
\end{enumerate}

\paragraph{Transverse $k$-point assembly.}
The function \texttt{\_assemble\_kpoint} performs the Fourier sum over
transverse periodic shifts.  For each $\vc{k}_\perp$ on a
Monkhorst--Pack mesh:
\begin{equation}\label{eq:assemble}
  D(\vc{k}_\perp)
  = \sum_{(n_x, n_y)} M(n_x, n_y)\;
    e^{2\pi i\,(k_x n_x + k_y n_y)},
\end{equation}
where $M(n_x, n_y)$ is the BTD sparse matrix for periodic shift
$(n_x, n_y)$ (already expanded in the transport direction).  The
$k$-points on the Monkhorst--Pack mesh are
\begin{equation}
  k_i = \frac{n_i + 0.5}{N_i} - 0.5 + \text{shift}_i,
  \quad n_i = 0, \ldots, N_i - 1.
\end{equation}
This is Convention~B from \cref{sec:dynmat}.  The result is a stack of BTD matrices
indexed by $(\omega, k_x, k_y)$.


% ======================================================================
\section{Ballistic Phonon}
\label{sec:harmonic}
% ======================================================================


% ----------------------------------------------------------------------
\subsection{Harmonic Force Constants}
\label{sec:fd}
% ----------------------------------------------------------------------

In the harmonic approximation, the crystal potential energy is expanded
to second order in the atomic displacements from equilibrium.  The
corresponding second-order coefficients are the harmonic force constants,
which quantify the restoring forces acting on the atoms. They are defined as
\begin{equation}\label{eq:fc}
  \Phi_{\alpha\beta}(0\kappa,\, l'\kappa')
  = \frac{\partial^2 E}
         {\partial u_\alpha(0\kappa)\,\partial u_\beta(l'\kappa')},
\end{equation}
where $u_\alpha(l\kappa)$ denotes the displacement of atom $\kappa$ in cell
$l$ along Cartesian direction $\alpha$.
$\Phi_{\alpha\beta}(0\kappa,l'\kappa')$ measures how a displacement of atom
$\kappa'$ in cell $l'$ along direction $\beta$ changes the force on atom
$\kappa$ in the home cell along direction $\alpha$. We have hence the
linear force--displacement relation
\begin{equation}\label{eq:force-disp}
  F_\alpha(0\kappa)
  = - \sum_{l'\kappa'\beta}
      \Phi_{\alpha\beta}(0\kappa,\, l'\kappa')\,
      u_\beta(l'\kappa').
\end{equation}
For a fixed pair of atoms $(0\kappa)$ and $(l'\kappa')$, the force constants
form a $3\times 3$ matrix, with rows corresponding to force
components and columns to displacement components.  If the unit cell
contains $N$ atoms, the full force-constant matrix has dimension
$(3N)\times(3N)$ for each lattice vector $\vc R_{l'}$.

The corresponding dynamical matrix is obtained from the mass-weighted
Fourier transform of $\Phi$, and its eigenvalues and eigenvectors yield
the phonon frequencies and polarization vectors (See \cite{dove_introduction_1993}).

In the finite-displacement method, the force constants are
obtained by displacing atoms in a supercell and computing the resulting DFT forces.
phonopy constructs a set of symmetry-inequivalent displaced
supercells, performs a self-consistent DFT calculation for each displaced
structure, and reconstructs the full force-constant tensor by inverting the
force--displacement relation \cref{eq:force-disp}.
Space-group symmetry is used to reduce
the number of independent displacements, and hence the number of required
DFT calculations (original \cite{parlinski_first-principles_1997}, phonopy \cite{togo_first_2015}).

The choice of supercell determines the real-space range over which the force
constants can be resolved before periodic wrap-around produces aliasing.
For a one-dimensional supercell of length $N$, the calculated force constants
can be written as
\begin{equation}
  \Phi_{\mathrm{per}}^{(N)}(r)
  = \sum_{p\in\mathbb Z} \Phi(r+pN).
\end{equation}
In a $2\times2\times2$ supercell, the resolved relative translations are limited
to those representable within that finite periodic cell. Contributions from more
distant cells are folded back onto these separations by periodic boundary
conditions. A larger supercell therefore allows force constants to be resolved
over a wider real-space range before aliasing occurs (see \cite[Chs.~2,~6]{dove_introduction_1993}).

% ----------------------------------------------------------------------
\subsection{Dynamical-Matrix}
\label{sec:dynmat}
% ----------------------------------------------------------------------

The dynamical matrix at wave vector $\vc{q}$ is the mass-weighted
Fourier transform of the force constants.  Phonopy and the quatrex
solver use different conventions:

\subsubsection{Convention A (phonopy)}

\begin{equation}\label{eq:DA}
  D_{\alpha\beta}^{(\mathrm{A})}(\kappa\kappa',\vc{q})
  = \frac{1}{\sqrt{m_\kappa m_{\kappa'}}}
    \sum_{l'} \Phi_{\alpha\beta}(0\kappa,\,l'\kappa')\;
    e^{2\pi i\,\vc{q}\cdot[\vc{R}_{l'} + \vc{\tau}_{\kappa'} - \vc{\tau}_\kappa]},
\end{equation}
where $\vc{q}$ is in fractional reciprocal coordinates (with the factor
$2\pi$ explicit).  The phase depends on the interatomic vector
$\vc{R}_{l'} + \vc{\tau}_{\kappa'} - \vc{\tau}_\kappa$ \cite{togo_first_2015}.

\subsubsection{Convention B (quatrex)}
\label{sec:convB}

quatrex assumes in its k-points code
\begin{equation}\label{eq:DB}
  D_{\alpha\beta}^{(\mathrm{B})}(\kappa\kappa',\vc{q})
  = \frac{1}{\sqrt{m_\kappa m_{\kappa'}}}
    \sum_{l'} \Phi_{\alpha\beta}(0\kappa,\,l'\kappa')\;
    e^{2\pi i\,\vc{q}\cdot \vc{R}_{l'}},
\end{equation}
where the phase involves only the lattice vector $\vc{R}_{l'}$.  This
makes $D^{(\mathrm{B})}$ a lattice Fourier series:
\begin{equation}\label{eq:fourier_series}
  D^{(\mathrm{B})}(\vc{q})
  = \sum_{\vc{R}} H(\vc{R})\;e^{2\pi i\,\vc{q}\cdot\vc{R}},
\end{equation}
where $H(\vc{R}) = \Phi(0,\vc{R}) / \sqrt{m\,m'}$ are real-valued
matrices.

\subsubsection{Transformation between conventions}
\label{sec:gauge}

The two conventions are related by a unitary similarity transformation
\begin{equation}\label{eq:gauge_eq}
  D^{(\mathrm{B})}(\vc{q})
  = P(\vc{q})\;D^{(\mathrm{A})}(\vc{q})\;P(\vc{q})^\dagger,
\end{equation}
with the diagonal gauge matrix
\begin{equation}\label{eq:P}
  P(\vc{q})_{\kappa\alpha,\,\kappa'\beta}
  = \delta_{\kappa\kappa'}\delta_{\alpha\beta}\;
    e^{2\pi i\,\vc{q}\cdot\vc{\tau}_\kappa},
\end{equation}
where $\vc{\tau}_\kappa$ is in fractional coordinates and each atomic
phase is repeated three times for the $x, y, z$ degrees of freedom.


% ----------------------------------------------------------------------
\subsection{Real-Space Block Extraction}
\label{sec:blocks}
% ----------------------------------------------------------------------

The inverse discrete Fourier transform (IDFT) of $D^{(\mathrm{B})}$
on a uniform $N_q^3$ mesh yields
\begin{equation}\label{eq:idft}
  H(n_x, n_y, n_z)
  = \frac{C}{N_q^3}
    \sum_{i,j,k=0}^{N_q - 1}
    D^{(\mathrm{B})}\!\left(\frac{i}{N_q}, \frac{j}{N_q}, \frac{k}{N_q}\right)
    e^{-2\pi i(n_x i + n_y j + n_z k)/N_q},
\end{equation}
where $C$ is the unit-conversion factor defined in
Section~\ref{sec:solver_units}.  Because $D^{(\mathrm{B})}$ is a
lattice Fourier series of real matrices, the blocks $H(\vc{R})$ are
real-valued.

For $N_q = 3$, the IDFT resolves lattice translations
$(n_x, n_y, n_z) \in \{-1, 0, +1\}^3$, giving up to 27 blocks with symmetry
\begin{equation}\label{eq:sym}
  H(\vc{R})^T = H(-\vc{R}),
\end{equation}
which implies Hermiticity of $D^{(\mathrm{B})}(\vc{q})$ at every $\vc{q}$.


% ----------------------------------------------------------------------
\subsection{Harmonic Phonon NEGF}
\label{sec:negf}
% ----------------------------------------------------------------------

For a given transverse wavevector $\vc{k}_\perp = (k_x, k_y)$, the
on-site and coupling blocks along the transport direction $z$ are
obtained by a partial Fourier transform: performing the 1D
IDFT gives:
\begin{align}
  H_{00}(\vc{k}_\perp)
    &= \frac{C}{N_{q_z}}
       \sum_{k=0}^{N_{q_z}-1}
       D^{(\mathrm{B})}\!\left(k_x, k_y, \frac{k}{N_{q_z}}\right),
  \label{eq:H00} \\
  H_{01}(\vc{k}_\perp)
    &= \frac{C}{N_{q_z}}
       \sum_{k=0}^{N_{q_z}-1}
       D^{(\mathrm{B})}\!\left(k_x, k_y, \frac{k}{N_{q_z}}\right)
       e^{-2\pi i k / N_{q_z}}.
  \label{eq:H01}
\end{align}
At non-zero $\vc{k}_\perp$, these blocks are in general
complex-valued: they are
\begin{equation}
    \sum_{n_x,n_y} H(n_x,n_y,n_z)\,e^{2\pi i(k_x n_x + k_y n_y)},
\end{equation}
a 2D Fourier sum of real coefficients which is complex whenever
$\vc{k}_\perp \neq 0$.  The coupling satisfies $H_{10} = H_{01}^\dagger$.

\subsubsection{Sancho--Rubio decimation}

The surface Green's function of a semi-infinite lead is obtained
iteratively.  Starting from bulk blocks $H_{00}$ and $H_{01}$:
\begin{enumerate}
  \item Initialise $\epsilon = \epsilon_s = (\omega^2 + i\eta)I -
        H_{00}$, $\alpha = -H_{10}$, $\beta = -H_{01}$.
  \item Iterate:
        $\epsilon_s \leftarrow \epsilon_s - \alpha\,\epsilon^{-1}\beta$,\quad
        $\epsilon   \leftarrow \epsilon - \alpha\,\epsilon^{-1}\beta
         - \beta\,\epsilon^{-1}\alpha$,\quad
        $\alpha     \leftarrow \alpha\,\epsilon^{-1}\alpha$,\quad
        $\beta      \leftarrow \beta\,\epsilon^{-1}\beta$.
  \item Converge when $\|\alpha\| + \|\beta\| < \text{tol}$.
  \item Surface Green's function: $g_s = \epsilon_s^{-1}$.
\end{enumerate}
For the left lead, $g_L = g_s(H_{00}, H_{01})$.  For the right lead,
$g_R = g_s(H_{00}, H_{01}^\dagger)$ (See original papers \cite{sancho_quick_1984, sancho_highly_1985}).

\subsubsection{Caroli (Landauer) formula}

\begin{align}
  \Sigma_L &= H_{10}\,g_L\,H_{01}, &
  \Sigma_R &= H_{01}\,g_R\,H_{10}, \\
  \Gamma_{L/R} &= i(\Sigma_{L/R} - \Sigma_{L/R}^\dagger), \\
  G^R &= \bigl[(\omega^2 + i\eta)\,I - H_D - \Sigma_L - \Sigma_R\bigr]^{-1}, &
  G^A &= (G^R)^\dagger.
\end{align}
The ballistic transmission at a single $\vc{k}_\perp$ is (Caroli)
\begin{equation}\label{eq:caroli}
  \mathcal{T}(\omega^2, \vc{k}_\perp)
  = \mathrm{Tr}\bigl[\Gamma_L\,G^R\,\Gamma_R\,G^A\bigr].
\end{equation}
The total transmission is averaged over a Monkhorst--Pack mesh of
transverse $k$-points:
\begin{equation}
  \bar{\mathcal{T}}(\omega^2)
  = \frac{1}{N_{k_\perp}} \sum_{\vc{k}_\perp}
    \mathcal{T}(\omega^2, \vc{k}_\perp).
\end{equation}
(E.g. reviews \cite{wang_nonequilibrium_2007, wang_quantum_2008} for formulas)

\subsubsection{quatrex Implementation}
\label{sec:phonon_solver}

The quatrex \texttt{PhononSolver} assembles the system matrix
\begin{equation}\label{eq:system}
  \bigl[G^R(\omega^2, \vc{k}_\perp)\bigr]^{-1}
  = (\omega^2 + i\eta)\,\mathbb{I}
    - D(\vc{k}_\perp) - \Sigma^R(\omega^2, \vc{k}_\perp),
\end{equation}
where $D(\vc{k}_\perp)$ is the BTD dynamical matrix loaded from the
\texttt{.mat} file.  For ballistic transport, self-energies come only from the open boundary conditions.

The lesser and greater self-energies at the contacts encode thermal
occupation via the Bose--Einstein distribution:
\begin{align}
  \Sigma^<_L &= i\,\Gamma_L\,n_\mathrm{BE}(\omega, T_L), &
  \Sigma^>_L &= i\,\Gamma_L\,\bigl[n_\mathrm{BE}(\omega, T_L) + 1\bigr],
\end{align}
and equivalently for the right contact at temperature $T_R$.  The
Keldysh Green's functions are then
\begin{align}
  G^< &= G^R\,\Sigma^<_\mathrm{total}\,(G^R)^\dagger, \\
  G^> &= G^R\,\Sigma^>_\mathrm{total}\,(G^R)^\dagger.
\end{align}

\subsubsection{Thermal conductance}

The ballistic thermal conductance per unit area is
\begin{equation}\label{eq:G_thermal}
  G(T) = \frac{1}{A_c}
      \int_0^\infty \frac{d\omega}{2\pi}\;
      \bar{\mathcal{T}}(\omega)\;\hbar\omega\;
      \frac{\partial n_\mathrm{BE}}{\partial T},
\end{equation}
where $A_c$ is the unit-cell cross-section perpendicular to transport,
$n_\mathrm{BE}(\omega, T)$ the Bose--Einstein distribution, and
\begin{equation}
  \frac{\partial n_\mathrm{BE}}{\partial T}
  = \frac{\hbar\omega}{k_B T^2}\;
    \frac{e^{\hbar\omega / k_B T}}{(e^{\hbar\omega / k_B T} - 1)^2}.
\end{equation}
The integral is evaluated numerically via the trapezoidal rule.


\subsection{Extracting a phonopy calculation to quatrex input}
\label{sec:extraction}

The block-tridiagonal Hamiltonian $H(\vc{R})$ (i.e. FC matrix) that quatrex requires
is reconstructed from the dynamical matrix $D(\vc{q})$ via an inverse discrete Fourier
transform and a gauge transformation then removes the basis-atom phases (should not be necessary,
could read the FC directly, but phonopy handles symmetries).  The pipeline is implemented in \texttt{phonon\_inputs}. The steps are:

% ======================================================================
\paragraph{Performing the DFT calculation.}
\label{sec:performing_dft}The force constants are computed from phonopy/QE:

\begin{enumerate}
  \item \textbf{Structure setup.}
        The unit cell is loaded via \texttt{structure.py:load\_structure}.
        A \texttt{Phonopy} object with a $2\times2\times2$ supercell is
        created by \\ \texttt{structure.py:create\_phonopy\_from\_config},
        which generates symmetry-inequivalent displacements.

  \item \textbf{DFT single-point calculations.}
        For each displaced supercell, a QE \texttt{pw.x} input file is
        written by \texttt{qe\_interface.py}.
        Existing converged outputs are detected and reused.
        Forces are parsed from QE output and converted from
        Ry/bohr to eV/\AA{} using the factor
        $1\;\text{Ry/bohr} = 25.711\;\text{eV/\AA}$.

  \item \textbf{Force constant extraction.}
        \texttt{force\_constants.py:produce\_force\_constants}
        passes the forces to phonopy, which inverts the
        force--displacement relation using space-group symmetry.

\end{enumerate}

For the Si reference case ($2\times2\times2$ supercell of the
8-atom conventional cell, 64 atoms), phonopy requires one
symmetry-inequivalent displacement.  The corresponding QE SCF
calculation takes $\sim$30\,s on 4 cores at \texttt{ecutwfc}~=~60\,Ry
with a $4\times4\times4$ $k$-mesh.

% ----------------------------------------------------------------------
\paragraph{Load the phonopy calculation and convert units.}
The phonopy calculation must be brought into eV/\AA{} units (See \texttt{structure.py:load\_phonopy\_calculation}).
For QE, phonopy stores lattice vectors in bohr and
force constants in Ry/bohr$^2$.  Hence we need to multiply (see \cref{sec:phonopy_units}):
\begin{itemize}
  \item Lattice: $a_\text{\AA} = a_\text{bohr} \times 0.52918$
  \item Force constants: $\Phi_{\text{eV/\AA\textsuperscript{2}}} = \Phi_{\text{Ry/bohr\textsuperscript{2}}} \times 48.587$
\end{itemize}

% ----------------------------------------------------------------------
\paragraph{Evaluate the dynamical matrix on a $\vc{q}$-mesh.}
The dynamical matrix $D^{(\mathrm{A})}(\vc{q})$ is evaluated on
an $N_{q_x} \times N_{q_y} \times N_{q_z}$ uniform mesh using
\texttt{phonon.get\_dynamical\_matrix\_at\_q(q)} in \texttt{convention.py:evaluate\_on\_mesh}.

The default mesh size is $N_q = 3$ in each direction, which resolves
lattice translations $\vc{R} \in \{-1,0,+1\}^3$ in the subsequent
IDFT, matching the range of a $2\times2\times2$ supercell.

% ----------------------------------------------------------------------
\paragraph{Gauge transform to Convention~B.}
At each $\vc{q}$, the phase convention is changed from~A (phonopy) to~B
(quatrex) (see \cref{sec:dynmat}):
\begin{equation}\label{eq:gauge}
  D^{(\mathrm{B})}(\vc{q})
  = P(\vc{q})\;D^{(\mathrm{A})}(\vc{q})\;P(\vc{q})^\dagger,
\end{equation}
with
$P(\vc{q}) = \mathrm{diag}(\ldots,
\underbrace{e^{2\pi i\,\vc{q}\cdot\vc{\tau}_\kappa},\,
            e^{2\pi i\,\vc{q}\cdot\vc{\tau}_\kappa},\,
            e^{2\pi i\,\vc{q}\cdot\vc{\tau}_\kappa}}_{x,y,z},
\ldots)$,
where $\vc{\tau}_\kappa$ are the fractional coordinates of basis atom
$\kappa$.

% ----------------------------------------------------------------------
\paragraph{Inverse DFT to real-space blocks.}
The Convention~B dynamical matrix is a pure lattice Fourier series
(see. \eqref{eq:fourier_series}), so the real-space blocks are by IDFT
\begin{equation}\label{eq:idft}
  H(n_x, n_y, n_z)
  = \frac{1}{N_{q_x} N_{q_y} N_{q_z}}
    \sum_{i,j,k}
    D^{(\mathrm{B})}\!\!\left(\frac{i}{N_{q_x}},\,
                               \frac{j}{N_{q_y}},\,
                               \frac{k}{N_{q_z}}\right)
    e^{-2\pi i\left(\frac{n_x i}{N_{q_x}}
                   + \frac{n_y j}{N_{q_y}}
                   + \frac{n_z k}{N_{q_z}}\right)}.
\end{equation}
The result $H(\vc{R})$ is real-valued because $D^{(\mathrm{B})}$ is a
lattice Fourier series of real force constants, the imaginary part
(order $10^{-15}$ relative) is discarded. The implementation is
in \texttt{convention.py:idft\_to\_blocks}.

The IDFT produces blocks in phonopy eigenvalue units,
eV/(amu$\cdot$\AA$^2$).  Each block is then multiplied by
\begin{equation}
  C = \frac{e\;[\text{J/eV}]}{m_u\;[\text{kg}] \times (10^{-10})^2\;[\text{m}^2/\text{\AA}^2]}
    \approx 9.648 \times 10^{27}\;\frac{\text{rad}^2}{\text{s}^2}
    \cdot \frac{\text{amu}\cdot\text{\AA}^2}{\text{eV}},
\end{equation}
converting them to $(\text{rad/s})^2$, and then reshaped to a dictionary
$\{(n_x,n_y,n_z) \mapsto H \in \mathbb{R}^{n_\text{dof}\times n_\text{dof}}\}$.

% ----------------------------------------------------------------------
\paragraph{Quatrex Inputs.}
The blocks are written to \texttt{dynamical\_matrix.mat} in\\
\texttt{quatrex\_writer.py:write\_dynamical\_matrix\_mat}
using \texttt{scipy.io.savemat} from the dict.
Additionally, we write (\texttt{quatrex\_writer.py:write\_all}):
\begin{itemize}
  \item \texttt{structure.xyz} --- extended XYZ with lattice vectors in
        the header and Cartesian atomic positions
        (\texttt{write\_structure\_xyz}).
  \item \texttt{quatrex\_config.toml} --- device configuration
        (transport direction, $k$-point grid, neighbour cutoff, OBC
        algorithm, etc.)\ \texttt{write\_quatrex\_config\_toml}.
\end{itemize}
For heterogeneous contacts, separate \texttt{.mat} files
for the left and right contacts are written when the corresponding
config fields are set.



% ======================================================================
\subsubsection{Validation}
\label{sec:extraction_validation}

All of the below is in \texttt{validation.py}:

% --- Acoustic sum rule ------------------------------------------------
\paragraph{Acoustic sum rule.}
(\texttt{check\_gamma\_point})
At the $\Gamma$ point \eqref{eq:fourier_series} reduces to
$D^{(\mathrm{B})}(\vc{0}) = \sum_{\vc{R}} H(\vc{R})$.
Translational invariance requires three eigenvalues to vanish
(the acoustic modes at $\Gamma$).  The check computes
\begin{equation}
  D_\Gamma = \sum_{\vc{R}} H(\vc{R}),
  \qquad
  \omega_n = \mathrm{sign}(\lambda_n)\,
             \sqrt{|\lambda_n|},
  \qquad
  \lambda_n = \text{eigenvalues of } D_\Gamma,
\end{equation}
and reports the three smallest $|\omega_n|$ in THz.
Also, $\max|D_\Gamma - D_\Gamma^T|$ should be approximately zero.

% --- Block symmetry ---------------------------------------------------
\paragraph{Block symmetry.}
(\texttt{check\_block\_symmetry})
From \eqref{eq:DB}
\begin{equation}\label{eq:block_sym}
  H(\vc{R})^T = H(-\vc{R})
  \qquad \forall\;\vc{R}.
\end{equation}
The check iterates over all block pairs $(\vc{R}, -\vc{R})$ in the
dictionary and reports $\max_{\vc{R}} \|H(\vc{R})^T - H(-\vc{R})\|_\infty$.


% --- Band structure comparison ----------------------------------------
\paragraph{Band structure comparison.}
(\texttt{compare\_band\_structure})
This check verifies that the extracted real-space blocks reproduce the
same phonon frequencies as phonopy's internal dynamical matrix.  At each
$\vc{q}$ on a path, two sets of eigenfrequencies are
computed:
\begin{enumerate}
  \item From phonopy:
        $D_\text{ph} = P(\vc{q})\,D^{(\mathrm{A})}(\vc{q})\,P(\vc{q})^\dagger$,
        converted to $(\text{rad/s})^2$ via $C$, then diagonalised.
  \item From the extracted blocks:
        $D_\text{rec}(\vc{q}) = \sum_{\vc{R}} H(\vc{R})\,
         e^{2\pi i\,\vc{q}\cdot\vc{R}}$,
        then diagonalised.
\end{enumerate}
The maximum eigenfrequency difference
$\max_{n,\vc{q}} |\omega_n^\text{ph} - \omega_n^\text{rec}|$ is
reported, a comparison plot is generated overlaying
the two band structures along the chosen path.

% ======================================================================
\section{Quatrex Modifications}
\label{sec:quatrex-mods}
% ======================================================================


% ----------------------------------------------------------------------
\subsection{Heterogeneous-Device Implementation}
\label{sec:heterogeneous}
% ----------------------------------------------------------------------

\subsubsection{Changes to quatrex for the L|D|R device}

The original unit-cell input mode constructs a homogeneous device: the
same dynamical matrix blocks are tiled across all transport cells.  To
support material interfaces (e.g.\ Si/Ge), the device must be
partitioned into three regions with distinct dynamical matrices:
\begin{equation}\label{eq:ldr}
  \underbrace{L \;\cdots\; L}_{n_L\;\text{cells}}
  \;\;\underbrace{D \;\cdots\; D}_{n_D\;\text{cells}}
  \;\;\underbrace{R \;\cdots\; R}_{n_R\;\text{cells}},
  \qquad n_L + n_D + n_R = n_\text{transport}.
\end{equation}

Four new fields were added to \texttt{DeviceConfig}
(\texttt{quatrex/core/config.py}):

\begin{center}
\begin{tabular}{lll}
  \toprule
  Field & Type & Purpose \\
  \midrule
  \texttt{num\_left\_cells}  & \texttt{PositiveInt | None} & $n_L$ \\
  \texttt{num\_right\_cells} & \texttt{PositiveInt | None} & $n_R$ \\
  \texttt{left\_matrix}      & \texttt{str | None} & left-lead \texttt{.mat} file \\
  \texttt{right\_matrix}     & \texttt{str | None} & right-lead \texttt{.mat} file \\
  \bottomrule
\end{tabular}
\end{center}

Necessarily $n_L + n_R < n_\text{transport}$ and it requires \texttt{construct\_from\_unit\_cell = true}.

\subsubsection{Matrix assembly}

Two new functions were added to \texttt{quatrex/device/inputs.py}.

\textbf{\texttt{\_expand\_heterogeneous\_matrix}} builds a BTD sparse
matrix for a single transverse periodic shift $(n_x, n_y)$.  Each
cell is assigned to one of the three regions according to its index.
The diagonal block always comes from the cell's own region:
\begin{equation}
  D_{ii} =
  \begin{cases}
    D_L & i < n_L, \\
    D_D & n_L \le i < n_L + n_D, \\
    D_R & i \ge n_L + n_D.
  \end{cases}
\end{equation}

\textbf{\texttt{\_create\_heterogeneous\_matrix}} loops over all
transverse periodic shifts, calling
\\\texttt{\_expand\_heterogeneous\_matrix} for each, and sums all shifts
at $\Gamma$ for the sparsity pattern.


\paragraph{Boundary coupling.}
The Hermiticity relation for real-space blocks at non-zero transverse
shift is
\begin{equation}\label{eq:hermiticity_perp}
  H(\vc{R}_\perp, n_z)^\dagger = H(-\vc{R}_\perp, -n_z),
\end{equation}
at each boundary the coupling block is taken from the region of the neighbouring
cell.  Specifically, at the left/device boundary ($i = n_L - 1$,
$i+1 = n_L$), the super-diagonal uses the left-lead super-diagonal
$H_{01}^L$ and the sub-diagonal uses $H_{10}^L$.  At the
device/right boundary both couplings come from the right-lead blocks
$H_{01}^R$ and $H_{10}^R$.


% ======================================================================
\section{Reference Calculations}
\label{sec:validation}
% ======================================================================

A naive Sancho--Rubio + Caroli python code is used as reference.


% ----------------------------------------------------------------------
\subsection{Bulk Validation}
\label{sec:bulk_validation}
% ----------------------------------------------------------------------

For a homogeneous Si system, the quatrex NEGF result and the
standalone Caroli reference are run with identical parameters and
compared directly.  Agreement between the two curves confirms that the
unit-cell assembly and ballistic phonon solver work properly.


% ----------------------------------------------------------------------
\subsection{Si/Ge Interface}
\label{sec:sige}
% ----------------------------------------------------------------------

\subsubsection{Model}

Following~\cite{guo_quantum_2020,latour_ab_2017,tian_enhancing_2012},
ballistic phonon transmission through a Si/Ge interface is computed
using a mass-mismatch model: the harmonic force constants are taken
from bulk Si throughout, and the only distinction between Si and Ge is
the atomic mass ($m_\text{Si} = 28.09$\,amu,
$m_\text{Ge} = 72.64$\,amu).  This approximation is justified because
Si and Ge have similar interatomic force constants but substantially
different masses.

The system is partitioned into three regions:
\begin{itemize}
  \item \textbf{Left lead}: Si force constants + Si masses (pure Si).
  \item \textbf{Device}: Si force constants + mixed Si/Ge masses
        (atoms with fractional coordinate $z < 0.5$ are Si; the rest
        are Ge).
  \item \textbf{Right lead}: Si force constants + Ge masses
        (mass-mismatch Ge).
\end{itemize}
The coupling between the left lead and the device is $H_{01}$ of the
Si lead, the coupling between the device and the right lead is
$H_{01}$ of the Ge lead.

\subsubsection{Computational details}

The Si force constants are obtained from a DFT calculation at the
equilibrium lattice constant $a = 5.466$\,\AA, using the conventional
8-atom diamond cell with a $2\times2\times2$ supercell (64 atoms), i.e.
nearest neighbour. DFT parameters: plane-wave cutoff 60~Ry, charge-density cutoff
480~Ry, $4\times4\times4$ $k$-point grid for the supercell, SCF
convergence threshold $10^{-10}$~Ry, Gaussian smearing 0.01~Ry, PBE
functional, \texttt{Si.pbe-n-rrkjus\_psl.1.0.0.UPF} pseudopotential.
Diamond symmetry reduces this to a single displacement.

The phonopy calculation was performed using phonopy's native QE
interface, which stores lengths in bohr and forces in Ry/bohr.  Force
constants are converted to eV/\AA$^2$ on loading
(Section~\ref{sec:phonopy_units}).  The Ge and interface phonopy
objects are created by cloning the Si object and substituting atomic
masses.


\begin{center}
\begin{tabular}{ll}
  \toprule
  Parameter & Value \\
  \midrule
  Transverse $k$-mesh & $20\times20$ \\
  Frequency range & $0.01$--$16.0$\,THz, 101 points \\
  Transport direction & $z$ (001) \\
  Total transport cells / $n_L$ / $n_D$ / $n_R$ & 5 / 2 / 1 / 2 \\
  Broadening $\eta$ & $\tfrac{1}{2}\Delta\omega^2 \approx 5.05\times10^{23}$ (rad/s)$^2$ \\
  $T_L$, $T_R$ & 301\,K, 299\,K \\
  OBC algorithm & Sancho--Rubio \\
  \bottomrule
\end{tabular}
\end{center}

\subsubsection{Results}

\paragraph{$\Gamma$-point frequencies.}
The Si calculation gives a maximum optical frequency of 15.55~THz at
$\Gamma$, consistent with reference values.  The
Ge clone (mass-mismatch) gives 9.67~THz, in exact agreement with the
mass scaling $f_\text{Ge} = f_\text{Si}\sqrt{m_\text{Si}/m_\text{Ge}}$.

\paragraph{Thermal conductance.}
The interfacial thermal conductance at 300~K is
\begin{equation}
  G = 196\;\text{MW/(m}^2\text{\,K)}.
\end{equation}

\begin{center}
\begin{tabular}{ll}
  \toprule
  Reference &  $G$ [MW/(m$^2$\,K)] \\
  \midrule
  Latour et al.\ (2017) \cite{latour_ab_2017} & 210 \\
  Tian et al.\ (2012) \cite{tian_enhancing_2012}     &209 \\
  Here                                &  196 \\
  \bottomrule
\end{tabular}
\end{center}
The $\sim$7\% difference from the  literature is
likely due to the coarser transverse $k$-mesh and DFT calculation.

\begin{figure}[htb]
  \centering
  \includegraphics[width=\textwidth]{fig/reference/si_ge_quatrex_comparison.png}
  \caption{Ballistic phonon transmission on a $20\times20$ transverse
    $k$-mesh with 101 frequency points.
    \textbf{Left:} Homogeneous bulk Si.
    \textbf{Right:} Si/Ge interface in the mass-mismatch approximation.}
  \label{fig:quatrex_comparison}
\end{figure}

% ======================================================================
\section{Anharmonic Phonon}
\label{sec:anharmonic}
% ======================================================================


% ----------------------------------------------------------------------
\subsection{Third-Order Force Constants}
\label{sec:fc3}
% ----------------------------------------------------------------------

The third-order (cubic) interatomic force constant is
\begin{equation}\label{eq:fc3_def}
  \Phi^{(3)}_{\alpha\beta\gamma}(0\kappa,\, l'\kappa',\, l''\kappa'')
  = \frac{\partial^3 E}
         {\partial u_\alpha(0\kappa)\;
          \partial u_\beta(l'\kappa')\;
          \partial u_\gamma(l''\kappa'')}.
\end{equation}
By commutativity of partial derivatives, $\Phi^{(3)}$ is fully
symmetric under any permutation of its three index triples
$(\alpha, 0\kappa)$, $(\beta, l'\kappa')$, $(\gamma, l''\kappa'')$.
Units: the raw FC3 from \texttt{thirdorder.py} is in eV/\AA$^3$, not mass-weighted.

\subsubsection{Computation with thirdorder.py}
\label{sec:thirdorder}

The FC3 is computed using \texttt{thirdorder\_espresso.py} via finite-difference procedure:
\begin{enumerate}
  \item \textbf{Sow}: generate displaced supercell configurations.  For
        each symmetry-inequivalent pair of atomic displacements, positive
        and negative perturbations are created.
  \item \textbf{DFT}: QE computes the forces for each
        displaced configuration.
  \item \textbf{Reap}: force differences yield the third derivatives
\end{enumerate}
The output file \texttt{FORCE\_CONSTANTS\_3RD} contains one block per
symmetry-reduced triplet $(0\kappa, l'\kappa', l''\kappa'')$, each
specifying the cell vectors $\vc{R}_{l'}$ and $\vc{R}_{l''}$ in
Cartesian coordinates (\AA), atom indices (1-based in the file,
converted to 0-based on read), and the $3\times3\times3$ tensor
$\Phi^{(3)}_{\alpha\beta\gamma}$ in eV/\AA$^3$.


% ----------------------------------------------------------------------
\subsection{Extracting FC3 Matrices}
\label{sec:fc3_extraction}
% ----------------------------------------------------------------------

\subsubsection{Remapping}

\texttt{thirdorder.py} uses the 2-atom FCC primitive cell, while the phonopy
harmonic force constants use the 8-atom
conventional cubic cell.  The FC3 must therefore be
remapped from the FCC basis to the conventional basis before being
combined with the harmonic dynamical matrix.

The FCC primitive vectors (in Cartesian coordinates) are
$\vc{a}_1 = \tfrac{a}{2}(0,1,1)$,
$\vc{a}_2 = \tfrac{a}{2}(1,0,1)$,
$\vc{a}_3 = \tfrac{a}{2}(1,1,0)$, with 2 atoms at fractional
positions $(0,0,0)$ and $(\tfrac{1}{4},\tfrac{1}{4},\tfrac{1}{4})$.
Each conventional-cell atom can be expressed as an FCC atom plus an
FCC lattice translation:
$\vc{r}_i^\mathrm{conv} + \vc{s}
= \vc{r}_{j}^\mathrm{FCC} + \sum_m n_m\,\vc{a}_m^\mathrm{FCC}$,
where $\vc{s}$ is an origin shift and $j \in \{0, 1\}$ is the FCC
sublattice index.

The phonopy conventional cell has an origin offset relative to the
standard diamond convention.  For silicon, the first
conventional-cell atom sits at fractional position
$(\tfrac{7}{8}, \tfrac{7}{8}, \tfrac{7}{8})$ rather than $(0,0,0)$.
The required Cartesian shift is
\begin{equation}\label{eq:shift}
  \vc{s} = \left(\frac{a_\mathrm{conv}}{8},\,
                  \frac{a_\mathrm{conv}}{8},\,
                  \frac{a_\mathrm{conv}}{8}\right).
\end{equation}

For each FC3 block with FCC atom indices
$(\kappa_i, \kappa_j, \kappa_k)$ and cell vectors
$(\vc{R}_j, \vc{R}_k)$:
\begin{enumerate}
  \item Map each FCC (atom, position) to the conventional cell: find
        the conventional atom index $i'$ and integer cell vector
        $\vc{n}$ such that
        $\vc{r}^\mathrm{FCC}_\kappa + \vc{R} - \vc{s}
         = \vc{r}^\mathrm{conv}_{i'} + \vc{n}\cdot\vc{A}_\mathrm{conv}$.
  \item Compute new cell vectors relative to atom $i$:
        $\vc{R}'_j = (\vc{n}_j - \vc{n}_i)\cdot\vc{A}_\mathrm{conv}$,
        $\vc{R}'_k = (\vc{n}_k - \vc{n}_i)\cdot\vc{A}_\mathrm{conv}$.
  \item Store the remapped block with conventional atom indices
        $(i', j', k')$ and new cell vectors $(\vc{R}'_j, \vc{R}'_k)$.
\end{enumerate}

For Si, there are 422 FC3 blocks (cutoff 3).

\subsubsection{Mass-weighting}

\begin{equation}\label{eq:fc3_mw}
  \Psi_{\alpha\beta\gamma}(\kappa\kappa'\kappa'')
  = \frac{\Phi^{(3)}_{\alpha\beta\gamma}(\kappa\kappa'\kappa'')}
         {\sqrt{M_\kappa\,M_{\kappa'}\,M_{\kappa''}}},
\end{equation}
with $M_\kappa$ in amu.  Units:
$[\Psi] = \mathrm{eV}/(\text{\AA}^3\cdot\mathrm{amu}^{3/2})$.

\subsubsection{Unit Conversion}

The harmonic conversion factor $C = e/(m_u A^2)$ converts
eV/(amu$\cdot$\AA$^2$) to (rad/s)$^2$.  The FC3 has one additional
factor of $1/(\text{\AA}\cdot\mathrm{amu}^{1/2})$ compared to FC2:
\begin{equation}\label{eq:conversion_fc3}
  \texttt{CONVERSION\_FC3}
  = \frac{C}{A\cdot\sqrt{m_u}}
  = \frac{e}{m_u^{3/2}\,A^3}
  \approx 2.368 \times 10^{51}
  \quad\bigl[\mathrm{J}/(\mathrm{m}^3\cdot\mathrm{kg}^{3/2})\bigr],
\end{equation}
converting eV/(\AA$^3$\,amu$^{3/2}$) to
J/(m$^3$\,kg$^{3/2}$), equivalently to
$(s^2\cdot\mathrm{m}\cdot\mathrm{kg}^{1/2})^{-1}$.


% ----------------------------------------------------------------------
\subsection{FC3 Assembly}
\label{sec:assembly_fc3}
% ----------------------------------------------------------------------

For a single NEGF transport cell the device region is one unit cell
in the transport direction, with a FC3 tensor
$\Psi_{abc}$ (a $n_\mathrm{dof}\times n_\mathrm{dof}\times n_\mathrm{dof}$
array, with $n_\mathrm{dof} = 3N_b = 24$ for Si).

\paragraph{Onsite approximation.}
Only FC3 blocks with both cell vectors equal to zero:
\begin{equation}
  \Psi_{abc}^\mathrm{onsite}
  =
    \frac{\Phi^{(3)}_{abc}(\vc{0}, \vc{0}, \vc{0})}
         {\sqrt{M_a M_b M_c}}.
\end{equation}

\paragraph{Full $\delta l = 0$ approximation (Approximation~III in \cite{guo_quantum_2020}).}
This approximation retains all FC3 blocks for which
both $\vc{R}'$ and $\vc{R}''$ lie in the same transport slab as the
origin, i.e.\ $\delta l' = \delta l'' = 0$, where
\begin{equation}
    \delta l' = \mathrm{round}\!\bigl([\vc{A}^{-T}\vc{R}']_z\bigr)
\end{equation}
so that
\begin{equation}\label{eq:fc3_full}
    \Psi_{abc}^\mathrm{full}
    = \sum_{\substack{\vc{R}',\vc{R}'' \\ \delta l'=0,\;\delta l''=0}}
      \frac{\Phi^{(3)}_{abc}(\vc{0}, \vc{R}', \vc{R}'')}
           {\sqrt{M_a M_b M_c}}.
\end{equation}

For Si, 193 of the 422 FC3 blocks satisfy $\delta l' = \delta l'' = 0$ (i.e. $R_z = 0$),
yielding 2126 non-zero elements out of $24^3 = 13\,824$ total.

% ----------------------------------------------------------------------
\subsection{Phonon--Phonon Self-Energy}
\label{sec:self_energy}
% ----------------------------------------------------------------------

\subsubsection{Self-energy formula)}

The phonon--phonon scattering self-energy for slab $L$ in the device,
at second order in the cubic anharmonicity (\cite[Eq. 8]{guo_quantum_2020}):
\begin{equation}\label{eq:sigma_guo}
  \Sigma^{<}_{ab}(\omega)
  = \frac{i\hbar}{2}\;
    \sum_{c,d,e,f} \Psi_{acd}\;
    \biggl[\int \frac{d\omega'}{2\pi}\;
      G^{<}_{ce}(\omega')\,G^{<}_{df}(\omega - \omega')
    \biggr]\;\Psi_{bef}
\end{equation}
where $G^<(\omega)$ is the lesser Green's function of the slab,
$\Psi$ is the mass-weighted FC3 in SI units, and $\hbar = 1.055\times
10^{-34}$\,J$\cdot$s.  The same formula with $G^<\to G^>$ gives
$\Sigma^>$.  In the absence of the Hilbert-transform contribution:
\begin{equation}\label{eq:sigma_r}
  \Sigma^R(\omega) \approx \tfrac{1}{2}
    \bigl[\Sigma^>(\omega) - \Sigma^<(\omega)\bigr].
\end{equation}

\subsubsection{Full Fourier-space self-energy (\cite[Eq. 14]{guo_quantum_2020})}
\label{sec:sigma_fourier}

For 3D nanostructures with transverse periodicity, the
self-energy involves a sum over all transverse wavevectors \cite[Eq. 14]{guo_quantum_2020}:
\begin{align}\label{eq:sigma_phph}
  \Sigma^{\gtrless}_{s,\,ll'_s}(&\omega;\,\vc{q}_\perp)
  = \frac{i\hbar}{2}
    \sum_{\substack{l_1 l_2 l_3 l_{4s} \\ j_1 j_2 j_3 j_4}}
    \frac{1}{N} \sum_{\vc{q}'_\perp}
    \tilde{\Phi}^{ij_1 j_2}_{l,\,l_{1s} l_{2s}}
      \!\left(\vc{q}'_\perp,\;\vc{q}_\perp - \vc{q}'_\perp\right)
    \notag \\
  &\qquad\times
    \tilde{\Phi}^{j_3 j_4}_{l',\,l_{3s} l_{4s}}
      \!\left(\vc{q}'_\perp - \vc{q}_\perp,\;-\vc{q}'_\perp\right)
    \int_{-\infty}^{\infty}\!\frac{d\omega'}{2\pi}
    \,G^{\gtrless,\,j_1 j_4}_{l_{1s} l_{4s}}(\omega';\,\vc{q}'_\perp)\,
    G^{\gtrless,\,j_2 j_3}_{l_{2s} l_{3s}}
      (\omega - \omega';\,\vc{q}_\perp - \vc{q}'_\perp),
\end{align}
where $l, l'$ index transport-direction slabs, $s$ denotes slab-local
atom indices, $i, j$ are Cartesian components, $N$ is the number of
transverse wavevectors, and the sum over $\vc{q}'_\perp$ enforces
transverse momentum conservation.

The quantity
$\tilde{\Phi}^{ij_1 j_2}_{l,\,l_{1s} l_{2s}}(\vc{q}_\perp, \vc{q}'_\perp)$
is the Fourier representation of the third-order FC matrix \cite[Eq. 16]{guo_quantum_2020}:
\begin{equation}\label{eq:fc3_fourier}
  \tilde{\Phi}^{ij_1 j_2}_{l,\,l_{1s} l_{2s}}
  (\vc{q}_\perp,\,\vc{q}'_\perp)
  = \sum_{\Delta\vc{R}_\perp}\sum_{\Delta\vc{R}'_\perp}
    \Phi^{ij_1 j_2}_{l l_{1} l_{2}}\;
    e^{-i\vc{q}_\perp\cdot\Delta\vc{R}_\perp}\;
    e^{-i\vc{q}'_\perp\cdot\Delta\vc{R}'_\perp},
\end{equation}
where
$\Delta\vc{R}_\perp = (l_y - l_{1y})\vc{a}_2 + (l_z - l_{1z})\vc{a}_3$
and
$\Delta\vc{R}'_\perp = (l_y - l_{2y})\vc{a}_2 + (l_z - l_{2z})\vc{a}_3$.

Unlike the harmonic FC matrix $\tilde{\Phi}^{(2)}(\vc{q}_\perp)$, which
depends on a single transverse wavevector, $\tilde{\Phi}^{(3)}$ depends
on two independent wavevectors.  Scaling hence:
\begin{center}
\begin{tabular}{lll}
  \toprule
  Object & Storage (dense) & Storage (NN-only) \\
  \midrule
  $\tilde{\Phi}^{(2)}(\vc{q}_\perp)$ &
    $N_k \times N_s^2$ &
    $N_k \times N_s \times z_2$ \\
  $\tilde{\Phi}^{(3)}(\vc{q}_\perp,\,\vc{q}'_\perp)$ &
    $N_k^2 \times N_s^3$ &
    $N_k^2 \times N_s \times z_3$ \\
  \bottomrule
\end{tabular}
\end{center}
where $N_s = 3N_\mathrm{atoms}$ is the number of DOFs per slab,
$z_2$ and $z_3$ are the average numbers of non-zero couplings per
site from harmonic and anharmonic interactions, and
$N_k = N_{q_y}\times N_{q_z}$ is the total number of transverse
wavevectors.

In the ballistic case each transverse wavevector $\vc{q}_\perp$ is
independent: the retarded Green's function decouples across
$\vc{q}_\perp$, so the problem is embarrassingly parallel over
$(\omega, \vc{q}_\perp)$.  In the anharmonic case, the sum
$N^{-1}\sum_{\vc{q}'_\perp}$ in Eq.~\eqref{eq:sigma_phph} couples
all transverse wavevectors: evaluating
$\Sigma(\omega; \vc{q}_\perp)$ requires
$G(\omega'; \vc{q}'_\perp)$ for every $\vc{q}'_\perp$.

% ----------------------------------------------------------------------
\subsection{Frequency Discretisation and FFT Convolution}
\label{sec:fft}
% ----------------------------------------------------------------------

\paragraph{Discrete frequency grid.}
$\omega_n = \omega_\mathrm{min} + n\,\Delta\omega$ for
$n = 0,\ldots,N_\omega-1$.  The broadening parameter is
$\eta = (\Delta\omega)^2 \times \eta_\mathrm{factor}$ with
$\eta_\mathrm{factor} = 0.5$ by default.

\paragraph{FFT convolution.}
The frequency integral in Eq.~\eqref{eq:sigma_guo} is discretised as
\begin{equation}
  C_{ce,df}[\omega_n]
  = \frac{\Delta\omega}{2\pi}\;
    \sum_{m} G^<_{ce}[\omega_m]\;G^<_{df}[\omega_n - \omega_m],
\end{equation}
computed via FFT: $C = \mathrm{IFFT}[\mathrm{FFT}(G_{ce})\cdot\mathrm{FFT}(G_{df})]$.

\paragraph{Full self-energy.}
\begin{equation}\label{eq:full-energy}
  \Sigma^<_{ab}[\omega_n]
  = \frac{i\hbar\,\Delta\omega}{4\pi}\;
    \sum_{c,d,e,f} \Psi_{acd}\;C_{ce,df}[\omega_n]\;\Psi_{bef}.
\end{equation}
This is implemented as a loop over $(c,d)$ and $(e,f)$ pairs with
sparsity checks on $|\Psi_{:,c,d}|$.  For Si
there are $\sim10^4$ active pairs.  The overall
computational cost is $\mathcal{O}(n_\mathrm{dof}^4 \times N_\omega
\log N_\omega)$.


% ----------------------------------------------------------------------
\subsection{SCBA Loop}
\label{sec:scba}
% ----------------------------------------------------------------------

\subsubsection{Algorithm}

\begin{enumerate}
  \item \textbf{Initialise}: $\Sigma^\mathrm{ph} = 0$.
  \item \textbf{Green's functions}: for each $\vc{q}_\perp$ and
        $\omega_n$, compute contact self-energies via Sancho--Rubio and
        solve:
        \begin{align}
          G^R &= \bigl[(\omega^2 + i\eta)\,\mathbb{I}
            - H_{00} - \Sigma^R_L - \Sigma^R_R
            - \Sigma^R_\mathrm{ph}\bigr]^{-1}, \\
          G^{<,>} &= G^R\,\Sigma^{<,>}_\mathrm{total}\,(G^R)^\dagger.
        \end{align}
  \item \textbf{Self-energy}: compute $\Sigma^{<,>,R}_\mathrm{ph}$
        from $G^{<,>}$ using Eq.~\eqref{eq:full-energy}.
  \item \textbf{Mixing}:
        $\Sigma^{(n+1)} = (1-\alpha)\,\Sigma^{(n)} +
        \alpha\,\Sigma^\mathrm{new}$, with $\alpha = 0.3$ by default.
  \item \textbf{Transmission}: Incorrect, change
  \item \textbf{Convergence check}:
        \begin{equation}
          \frac{\sum_n |\mathcal{T}^{(k)}_n - \mathcal{T}^{(k-1)}_n|}
               {\sum_n |\mathcal{T}^{(k-1)}_n|}
          < \mathrm{tol}.
        \end{equation}
  \item Repeat from step~2 until converged or maximum iterations
        reached.
\end{enumerate}

\subsubsection{Computational cost per iteration}

Per SCBA iteration:
\begin{itemize}
  \item \textbf{Green's functions}: $N_\perp \times N_\omega$ matrix
        inversions of size $n_\mathrm{dof}\times n_\mathrm{dof}$, plus
        Sancho--Rubio evaluations.
  \item \textbf{Self-energy}:
        $\mathcal{O}(n_\mathrm{dof}^4 \times N_\omega\log N_\omega)$
        from the FFT convolution loop.
  \item \textbf{Transmission}: TODO
\end{itemize}
The self-energy step dominates for large $n_\mathrm{dof}$.


% ----------------------------------------------------------------------
\subsection{Anharmonic Results and Validation}
\label{sec:anharmonic_results}
% ----------------------------------------------------------------------

\subsubsection{Test configuration}

Silicon, 8-atom conventional cell, $a = 5.4662$\,\AA; harmonic FC
from phonopy (QE, $2\times2\times2$ supercell); FC3 from
\texttt{thirdorder.py} (2-atom FCC cell, 422 blocks), remapped to the
conventional cell.  Transport: $z$-direction, 300\,K, 51 frequency
points from 0.5 to 15.0\,THz, mixing $\alpha = 0.3$, FC3 only keeping
transport slab contribution (see \cref{sec:assembly_fc3})

\begin{table}[h]
\centering
\caption{Self-energy magnitude and thermal conductance as a function
of the transverse $q$-mesh.}
\label{tab:scaling_anh}
\begin{tabular}{cccccc}
  \toprule
  $q$-mesh & $N_\perp$ &
    $|\Sigma^R|_\mathrm{max}/|H_{00}|_\mathrm{max}$ &
    SCBA iter &
    $G_\mathrm{ball}$ &
    $G_\mathrm{anh}$ \\
  & & & & [MW/(m$^2$K)] & [MW/(m$^2$K)] \\
  \midrule
  $2\times2$ & 4  & 0.49 & 11 & 1096 &  913 \\
  $4\times4$ & 16 & 0.20 &  8 & 1162 & 1005 \\
  \bottomrule
\end{tabular}
\end{table}


\subsubsection{Transmission}

\begin{figure}[h]
  \centering
  \includegraphics[width=\textwidth]{fig/reference/anharmonic_test.png}
  \caption{Left: ballistic (blue) and anharmonic SCBA (red dashed)
    phonon transmission for bulk Si at 300\,K ($4\times4$ $q$-mesh,
    51 frequency points).
    Right: transmission ratio
    $\mathcal{T}_\mathrm{anh}/\mathcal{T}_\mathrm{ball}$. (Wrong cause wrong formula)}
  \label{fig:transmission}
\end{figure}

Acoustic modes ($<$6~THz): transmission ratio $\approx 1.0$;
few available decay channels because conservation of energy limits
the phase space for $\omega\to\omega'+\omega''$.  Optical modes
(10--15~THz): ratio $\approx 0.7$--$0.9$; higher phonon density of
states provides more decay channels.

% ======================================================================
\section{Solver Architecture and Scaling Analysis}
\label{sec:scaling_arch}
% ======================================================================

Notation follows~\cite{vetsch_ab-initio_2025}:

\begin{table}[htb]
\centering
\begin{tabular}{cl}
  \toprule
  Symbol & Description \\
  \midrule
  $N_E$                & Number of energy / frequency grid points \\
  $N_B$                & Number of BTD blocks \\
  $\widetilde{N}_{BS}$ & Size of one BTD block (DOFs per transport cell) \\
  $N_{AO}$             & Total matrix dimension ($= N_B\times\widetilde{N}_{BS}$) \\
  $N_k$                & Number of transverse $k$-points ($=\prod_i k_i$ for $k_i>1$) \\
  $r_\mathrm{cut}$     & Interaction cutoff radius (determines sparsity) \\
  $N_\mathrm{nnz}$     & Non-zero elements per matrix \\
  $P_S$                & Number of spatial domain partitions (nested dissection) \\
  $P$                  & Total MPI ranks (GPUs) \\
  $N_\mathrm{iter}$    & Number of SCBA iterations \\
  \bottomrule
\end{tabular}
\caption{Notation.}
\label{tab:notation}
\end{table}


% ----------------------------------------------------------------------
\subsection{Data Distribution Model}
\label{sec:data_distribution}
% ----------------------------------------------------------------------

quatrex uses a two-level MPI communicator hierarchy
(\texttt{QuatrexCommunicator} in \\\texttt{qttools/comm/comm.py}):
\begin{itemize}
  \item \textbf{Stack communicator} (\texttt{comm.stack}): distributes
        the energy$\times k$-point grid across ranks.  Each rank owns a
        contiguous slice of $N_E\times N_k$ stack entries; operations at
        different energies and $k$-points are embarrassingly parallel.
  \item \textbf{Block communicator} (\texttt{comm.block}): distributes
        the spatial BTD blocks across ranks within each stack group,
        enabling the distributed RGF algorithm (nested-dissection
        spatial domain decomposition) of~\cite{vetsch_ab-initio_2025}.
\end{itemize}

\paragraph{DSDBSparse data structure.}
All Green's functions, self-energies, dynamical matrices, and screened
interactions are stored as \texttt{DSDBSparse} (Distributed Stack of
Distributed Block-accessible Sparse) matrices.  Each matrix has shape
\\$(n_{E,\text{local}}, n_{k_1},\ldots)\times(N_{AO}\times N_{AO})$,
but only the non-zero entries within the BTD band are stored:
$N_\mathrm{nnz} = \mathcal{O}(r_\mathrm{cut}\,\widetilde{N}_{BS}\,N_{AO})$.

A DSDBSparse matrix exists in one of two states, toggled by
\texttt{dtranspose()}:
\begin{enumerate}
  \item \textbf{Stack-distributed} (default): each rank holds all
        $N_\mathrm{nnz}$ entries for its local stack slice.  Used
        during OBC computation and RGF.
  \item \textbf{NNZ-distributed}: each rank holds all $N_E N_k$ stack
        entries for its local subset of $(i,j)$ non-zero pairs.  Used
        during energy convolutions.
\end{enumerate}
When symmetry is enabled (\texttt{symmetry=True}), only the upper
triangle is stored, halving both memory and communication.  The
transposition all-to-all costs
$\mathcal{O}(N_E N_k\,r_\mathrm{cut}\,\widetilde{N}_{BS}\,N_{AO})$
bytes.


% ----------------------------------------------------------------------
\subsection{Component-Wise Scaling}
\label{sec:component_scaling}
% ----------------------------------------------------------------------

\subsubsection{System matrix assembly}

For each local stack entry the system matrix is assembled as
\begin{equation}
  \widetilde{M}(\omega^2) = (\omega^2 + i\eta)\,\mathbb{I} - D - \Sigma^R_\mathrm{scatt},
\end{equation}
involving element-wise subtraction on the BTD, touching
$\mathcal{O}(N_B\,\widetilde{N}_{BS}^2)$ elements per stack entry.

\subsubsection{Open boundary conditions}

OBC is computed at both contacts for each stack entry.

\paragraph{Sancho--Rubio}: iterations involving three
$\widetilde{N}_{BS}\times\widetilde{N}_{BS}$ matrix multiplications
and one inversion.  Computation:
$\mathcal{O}(N_E N_k\,N_\mathrm{SR}\,\widetilde{N}_{BS}^3)$; memory:
$\mathcal{O}(N_E N_k\,\widetilde{N}_{BS}^2)$; no communication (each
contact computed only on the boundary rank).

\paragraph{Spectral / Beyn}: solves a polynomial
eigenvalue problem of size $N_U\widetilde{N}_{BS}$ (where $N_U$ is the
number of block sections) via a contour integral. \\ Computation:
$\mathcal{O}(N_E N_k\,N_\mathrm{quad}\,(N_U\widetilde{N}_{BS})^3)$;
memory: $\mathcal{O}(N_E N_k\,(N_U\widetilde{N}_{BS})^2)$; no
communication.

\subsubsection{RGF selected solve}

The forward sweep computes $N_B$ Schur complements; the backward
sweep computes $N_B$ matrix multiplications.  Each involves
$\widetilde{N}_{BS}\times\widetilde{N}_{BS}$ dense matrices.  For
$P_S = 1$: computation
$\mathcal{O}(N_E N_k\,N_B\,\widetilde{N}_{BS}^3)$.  For distributed
RGF ($P_S > 1$), the nested-dissection scheme splits the BTD matrix
into $P_S$ spatial partitions; a reduced system of size
$P_S\,\widetilde{N}_{BS}$ must be solved between the forward and
backward passes:
\begin{equation}
  \text{Computation: }
  \mathcal{O}\!\left(N_E N_k\!\left(\frac{N_B}{P_S} + P_S\right)\widetilde{N}_{BS}^3\right).
\end{equation}
Communication cost is $\mathcal{O}(N_E(P_S-1)\widetilde{N}_{BS}^2)$
(reduced system only).

\subsubsection{Energy convolutions}

For GW, the scattering self-energies and polarization are computed via FFT
energy convolutions.  In the nnz-distributed state, the full energy
spectrum for each $(i,j)$ non-zero entry is local, enabling an
in-place FFT:
\begin{equation}
  \text{Computation: }
  \mathcal{O}\!\left(N_E\log N_E \cdot N_k \cdot \frac{r_\mathrm{cut}\,
  \widetilde{N}_{BS}\,N_{AO}}{P}\right).
\end{equation}
When transverse $k$-points are present the FFT is multi-dimensional
(over both energy and $k$-axes), costing
$\mathcal{O}(N_E N_k\log(N_E N_k)\cdot N_\mathrm{nnz,local})$.

The retarded self-energy $\Sigma^R$ is obtained from $\Sigma^>-\Sigma^<$
via a Hilbert transform, implemented as additional FFT convolutions with
a $1/E$ kernel, doubling the FFT cost for each retarded component.

On GPUs, memory for FFT buffers is dynamically batched: the code
estimates available device memory, computes a batch size, and processes
$(i,j)$ entries in batches to avoid OOM.


\subsubsection{Summary table}

\begin{table}[htb]
\centering
\footnotesize
\begin{tabular}{lllc}
  \toprule
  Component & Computation & Memory (per rank) & Comm. \\
  \midrule
  OBC (Sancho--Rubio) &
    $N_E N_k\,\widetilde{N}_{BS}^3$ &
    $N_E N_k\,\widetilde{N}_{BS}^2$ & --- \\
  G-assembly &
    $N_E N_k\,N_B\,\widetilde{N}_{BS}^2$ &
    $N_E N_k\,N_\mathrm{nnz}$ & --- \\
  RGF ($P_S=1$) &
    $N_E N_k\,N_B\,\widetilde{N}_{BS}^3$ &
    $N_E N_k\,N_B\,\widetilde{N}_{BS}^2$ & --- \\
  RGF ($P_S>1$) &
    $N_E N_k(N_B/P_S+P_S)\widetilde{N}_{BS}^3$ &
    $N_E N_k(N_B+P_S)\widetilde{N}_{BS}^2$ &
    $N_E P_S\widetilde{N}_{BS}^2$ \\
  W-assembly &
    $N_E N_k\,N_{AO}\,P_{BW}^2$ &
    $N_E N_k\,N_\mathrm{nnz}$ & halo \\
  FFT convolution &
    $N_E N_k\,N_\mathrm{nnz}\,\log N_E$ &
    $N_E N_k\,N_\mathrm{nnz}$ & --- \\
  All-to-all transpose &
    $N_E N_k\,N_\mathrm{nnz}$ &
    $N_E N_k\,N_\mathrm{nnz}$ &
    $N_E N_k\,N_\mathrm{nnz}$ \\
  \bottomrule
\end{tabular}
\caption{Scaling per SCBA iteration.  Prefactors and $P$-dependent
  divisors are omitted.  Memory in units of $16\,\text{bytes}\times$
  count (complex128).  The dominant terms for large devices are the
  RGF (cubic in block size) and the all-to-all transposition
  (communication-bound).}
\label{tab:scaling_summary}
\end{table}


% ----------------------------------------------------------------------
\subsection{Memory Scaling}
\label{sec:memory}
% ----------------------------------------------------------------------

Each DSDBSparse matrix stores $N_\mathrm{nnz}$ complex-128 values per
stack entry:
$\text{Memory per matrix} = 16\,n_\mathrm{stack}\,N_\mathrm{nnz}$~bytes.

\begin{center}
\begin{tabular}{lcc}
  \toprule
  Subsystem & Matrices & Count \\
  \midrule
  Electron (G-solver) &
    $G^<,G^>,G^R$; $\Sigma^{<,>,R}$;
    $\Sigma^{<,>,R}_\mathrm{prev}$; system matrix; $D$
    & 11 \\
  Coulomb screening (W-solver) &
    $P^<,P^>,P^R$; $W^{<,>}$; system; $L^{<,>}$; $V$
    & 9 \\
  Phonon ballistic (D-solver) &
    $G^<,G^>,G^R$; system matrix; $D$
    & 5 \\
  Combined peak (NEGF+scGW) & $\sim$14 & \\
  \bottomrule
\end{tabular}
\end{center}

% ----------------------------------------------------------------------
\subsection{Anharmonic Scaling and Resource Estimates}
\label{sec:phph_scaling}
% ----------------------------------------------------------------------

\subsubsection{Computational complexity of the self-energy}

Breaking down Eq.~\eqref{eq:sigma_phph}:

\paragraph{Frequency convolution (inner loop).}
The $\int d\omega'$ integral is an energy convolution of two Green's
functions, $G(\omega')\otimes G(\omega-\omega')$.  This is evaluated
via FFT in $\mathcal{O}(N_\omega\log N_\omega)$ per $(i,j)$ entry per
$\vc{q}'_\perp$.

\paragraph{$\vc{q}'_\perp$ summation (outer loop).}
For each target $\vc{q}_\perp$, the sum over $\vc{q}'_\perp$ involves
$N_k$ terms.  Because the two Green's function arguments are
$\vc{q}'_\perp$ and $\vc{q}_\perp-\vc{q}'_\perp$, this is a
transverse-momentum convolution:
\begin{equation}
  \sum_{\vc{q}'_\perp}
  X(\vc{q}'_\perp)\,Y(\vc{q}_\perp - \vc{q}'_\perp)
  = \mathrm{IFFT}\bigl[\mathrm{FFT}[X]\cdot\mathrm{FFT}[Y]\bigr](\vc{q}_\perp).
\end{equation}
The multi-dimensional FFT already implemented in quatrex's
\texttt{fft\_convolve\_kpoints} takes the FFT jointly over
$(\omega, q_y, q_z)$ axes, converting both the energy convolution
\emph{and} the $\vc{q}$-summation into
$\mathcal{O}(N_\omega N_k\log(N_\omega N_k))$ per $(i,j)$ entry.

\paragraph{FC3 weighting.}
The FC3 tensor $\tilde{\Phi}^{(3)}$ appears as a pre- and post-factor
that contracts internal indices $(j_1,j_2,l_1,l_2)$ with the product
of Green's functions.  In Approximation~(III) of \cite{guo_quantum_2020} (diagonal blocks
only) the contraction involves $\mathcal{O}(N_s^2)$ index pairs per
$(l,l')$ slab pair.

\paragraph{Total operation count.}
\begin{equation}\label{eq:phph_cost}
  C_\Pi = \mathcal{O}\!\left(
    N_\mathrm{nnz} \times N_\omega\,N_k\,\log(N_\omega N_k)
  \right),
\end{equation}
where $N_\mathrm{nnz}^{(\Pi)} \approx N_d\times N_s^2$ (diagonal
blocks only, Approximation~III) or $\approx 3N_d N_s^2$ (full BTD).
\cite{guo_quantum_2020} show that the full BTD structure is needed for quantitative
accuracy (8--12\% error with diagonal-only; divergence with the even
simpler $3\times3$ diagonal approximation).

\subsubsection{Comparison with GW polarization}

The phonon--phonon self-energy has the same mathematical structure as
the GW polarization $P = G\cdot G$, and in quatrex's
\texttt{PCoulombScreening.compute()} an identical FFT-based convolution
is already performed.  Differences:
\begin{enumerate}
  \item \textbf{FC3 weighting}: the phonon self-energy is a
        ``sandwiched'' convolution
        $\Phi^{(3)}\cdot(G\otimes G)\cdot\Phi^{(3)}$, requiring
        contraction with the anharmonic tensor before and after the FFT.
  \item \textbf{Grid size}: the phonon frequency grid
        ($N_\omega\sim$100--200) is smaller than typical electronic
        energy grids ($N_E\sim$1000--5000), but the $q$-mesh is denser
        ($6\times6$ to $20\times20$ vs.\ often $1\times1$ or $3\times3$
        for electrons).
  \item \textbf{Block size}: phonon DOFs are 3 per atom (vs.\
        $\mathcal{O}(10)$ orbitals/atom for electrons), giving smaller
        $\widetilde{N}_{BS}$ but the same $N_B$ scaling.
\end{enumerate}

For ballistic transport, doubling the $q$-mesh from $n\times n$ to
$2n\times 2n$ increases $N_k$ by $4\times$ and the total cost scales
as $\mathcal{O}(N_k)$. For anharmonic transport the self-energy's $q$-convolution
couples all wavevectors, so cost scales as $\mathcal{O}(N_k\log N_k)$
per entry \emph{and} the FC3 storage grows as $\mathcal{O}(N_k^2)$
(from the second wavevector dependence).  Concretely, the FC3 storage
per rank is
\begin{equation}\label{eq:fc3_mem}
  M_{\Phi^{(3)}} \approx N_k \times N_\mathrm{trip} \times 16\;\text{bytes},
\end{equation}
where $N_\mathrm{trip}$ is the number of non-zero triplet elements per
wavevector.

\subsubsection{Resource estimates}

\paragraph{Si thin film (\cite{guo_quantum_2020}).}
8 atoms/cell ($N_s = 24$), $N_d = 20$, $6\times6$ mesh
($N_k = 36$), $N_\omega = 101$.  Stack size: 3\,636.
$N_\mathrm{nnz} \approx 20\times576\times3 = 34\,560$.
Memory per matrix: $16\times3\,636\times34\,560 = 2.0$\,GB.
Broken down by type:
\begin{itemize}
  \item $G^<, G^>, G^R$ (3 matrices): $3\times2.0 = 6.0$\,GB.
  \item $\Sigma^<, \Sigma^>, \Sigma^R$ + three previous-iteration copies
        (6 matrices): $6\times2.0 = 12.0$\,GB.
  \item System matrix + dynamical matrix (2 matrices): $2\times2.0 =
        4.0$\,GB.
  \item Total DSDBSparse (11 matrices): $\sim$\textbf{22\,GB}.
  \item FFT buffers (padded to $2N_\omega\times N_k$,
        $\sim1.5\times$ above): $\sim$\textbf{5\,GB} temporary.
  \item Total $\sim$\textbf{27\,GB per rank}.
\end{itemize}

Same geometry, $N_k = 400$.  Stack size: 40\,400.
Memory per matrix: 22.4\,GB; 11 matrices: $\sim$246\,GB
